% Source-derived chapter 01: Introduction
% Original source: geometric-bogomolov-conjecture-arbitrary-characteristics
\chapter{Introduction}
\label{geometric-bogomolov-conjecture-arbitrary-characteristics-chapter-01}
\source{geometric-bogomolov-conjecture-arbitrary-characteristics}

\begin{remark}[Source section]
\label{geometric-bogomolov-conjecture-arbitrary-characteristics-chapter-01-source}
\source{geometric-bogomolov-conjecture-arbitrary-characteristics}
\source{geometric-bogomolov-conjecture-arbitrary-characteristics}
This chapter reproduces the corresponding section of the retrieved
LaTeX source, including its definitions, statements, examples, and
proofs. It has not yet been translated into Lean.
\notready
\end{remark}

% The source section title was: Introduction
The goal of this paper is to prove the full geometric Bogomolov conjecture.
We first reduce it to the case that the extension of the base fields has transcendence degree 1, and then we prove the later case by intersection theory in algebraic geometry. 
The proof uses Yamaki's reduction theorem on the geometric Bogomolov conjecture and the Manin--Mumford conjecture proved by Raynaud and Hrushovski.  
%\subsection{The geometric Bogomolov conjecture}
%%%%%%%%%%%%%%%%%%%%%%%%%%%%%%%%%%%%%%%%%%%%%%

%%%%%%%%%%%%%%%%%%%%%%%%%%%%%%%%%%%%%%%%%%%%%%
\subsection{Abelian varieties and heights}\label{par:1.1.1}
\source{geometric-bogomolov-conjecture-arbitrary-characteristics}
%%%%%%%%%%%%%%%%%%%%%%%%%%%%%%%%%%%%%%%%%%%%%%
 
Let $k$ be an algebraically closed field.
%Let $k$ be a field.
Let $K/k$ be a finitely generated field extension of transcendence degree $\trdeg(K/k)\geq 1$.
Let $A$ be an abelian variety over $K$ of dimension $g$. 
Let $L$ be a symmetric and ample line bundle over $A$.
To define canonical height of subvarieties of $A$, we need choose integral models. 

\subsubsection{Integral models} \label{sectioncanheight}
\source{geometric-bogomolov-conjecture-arbitrary-characteristics}

A \emph{projective model of $K/k$} is a normal projective variety $S$ over $k$ with function field $K$.
It follows that $\dim S=\trdeg(K/k)$.

A \emph{polarization of $K/k$} is a pair $(S, \sM)$ consisting of a projective model $S$ of $K/k$ and an ample line bundle $\sM$ over $S$. 

An \emph{integral model of $(A,L)$} over $S$ is a pair $(\sA, \pi, \sL)$ where: 
\begin{points}
\item[$\bullet$] $\sA$ is a projective variety over $k$;
\item[$\bullet$] $\pi: \sA\to S$ is a projective morphism whose generic fiber is $A$;
\item[$\bullet$] $\sL$ is a line bundle on $\sA$ extending $L$.
\end{points}
We usually abbreviate $(\sA, \pi, \sL)$ as $(\sA, \sL)$.
%One may construct a model $(\sA, m, \sL)$ of $(A,L)$ such that $\sL$ is ample over $\sA$ as follows. 
%There is $m\geq 1$ such that $L^{\otimes m}$ is very ample. We use $L^{\otimes m}$ to embed $A$ into
%$\P_{K}^N$ for some $N>0$. The Zariski closure $\sA$ of $A$ inside $\P^N \times
%S$ is an irreducible projective variety. We write
%$\pi:\sA\rightarrow S$ for the projection. The pullback $\sL'$ of $\sO(1)$ on
%$\P^N\times S$ to $\sA$ is $\pi$-ample and nef. Set $\sL = \sL'\otimes \pi^*\sM,$ which is ample. 
% The restriction of $\sL$ to $A$ is isomorphic to $L^{\otimes m}$.  
% Finally, replacing  $\sA$ by its normalization, we assume that $\sA$ is normal 
% ($\sL$ remains ample on the normalization). 

\subsubsection{Canonical heights}

Let $X$ be any closed subvariety of $A$ over $K$.
The {\bf{naive height}} of $X$ with respect to the polarization $(S,\sM)$
and the integral model $(\sA, \sL)$ is defined as 
\begin{equation}
\label{eq:naiveheight}
\source{geometric-bogomolov-conjecture-arbitrary-characteristics}
h_{(\sA,\sL)}^{\sM}(X):=\frac{\sL^{\dim X+1}\cdot (\pi^*\sM)^{\dim S-1}\cdot \sX}{(1+\dim X)\deg_L(X)},
\end{equation}
where $\sX$ is the Zariski closure of $X$ in $\sA$, the numerator is the intersection number in $\sA$, and $\deg_L(X)=L^{\dim X}\cdot X$
is the intersection number in $A$.

By Tate's limiting argument, the {\bf{canonical height}} of $X$ with respect to the polarization $(S,\sM)$
and the line bundle $L$ over $A$ is the limit
\begin{equation}\label{eq:ntlimit}
\source{geometric-bogomolov-conjecture-arbitrary-characteristics}
\hat{h}_L^{\sM}(X)
:=\lim_{n\to +\infty} \frac{1}{n^2}h_{(\sA,\sL)}([n]X).
\end{equation}
Here $[n]X$ is the image of $X$ under the multiplication map $[n]:A\to A$.
The limit exists, depends on $(S,\sM)$ and $L$, but is independent of the integral model $(\sA, m, \sL)$; see Gubler's work \cite[Theorem 3.6]{Gubler2007} and \cite[Theorem 11.18]{Gubler:Pisa03}.
Moreover, $\hat{h}_L^{\sM}(X)\geq0$ by choosing $\sL$ to be ample in the definition.
Note that our heights are normalized by the degree and dimension in the denominator. 

The definition of the heights can be easily extended to subvarieties of $A_{\overline{K}}=A\otimes_K\overline{K}$. In fact, let $X'$ be a closed subvariety of $A_{\overline{K}}$.
Let $X$ be the minimal closed subvariety of $A$ such that $X_{\overline K}$ contains $X'$ in $A_{\overline{K}}$, or equivalently $X$ is the image of the composition 
$X'\to A_{\overline{K}}\to A$ as schemes. 
Then we simply define
\begin{equation}\label{eq:ntlimit}
\source{geometric-bogomolov-conjecture-arbitrary-characteristics}
h_{(\sA,m,\sL)}^{\sM}(X')
:=h_{(\sA,m,\sL)}^{\sM}(X),\qquad
\hat{h}_L^{\sM}(X')
:=\hat{h}_L^{\sM}(X).
\end{equation}
For simplicity, we usually omit the dependence on $L$ and $(S,\sM)$ and write 
$$\hat{h}(X)=\hat{h}_L^{\sM}(X),\qquad
\hat{h}(X')=\hat{h}_L^{\sM}(X').$$

In the special case of $\dim X'=0$, we get the canonical height  for points 
$$\hat{h}=\hat{h}_L^{\sM}:A(\overline{K})\to [0,+\infty)$$ 
with respect to $L$ and $(S,\sM)$.

In the special case of $\trdeg(K/k)=1$, the situation is much easier, since the the projective model $S$ is unique and the line bundle $\sM$ is not used in the definition of the heights.


\subsubsection{Small points and special subvarieties}
For any subvariety $X$ of  $A_{\overline{K}}$ and any $\epsilon> 0$, we set 
\begin{equation}
X(\epsilon):=\{x\in X(\overline{K})|\,\, \hat{h}(x)<\epsilon\}.
\end{equation}
We say that $X$ \emph{contains a dense set of small points of $A/K/k$} if $X(\epsilon)$ is Zariski dense in $X$ for all $\epsilon >0$.
This notion is actually independent of the choice of $L$ and $(S,\sM)$ to define the canonical height $\hat h$, since any two such heights bound each other up to positive  multiples.

Let $(A^{\overline{K}/k}, \tr)$ be the $\overline{K}/k$-trace of  $A_{\overline{K}}$; it
is the final object of the category of pairs $(C,f)$, where $C$ is an abelian variety over $ k$ and $f$ is 
a morphism from $C\otimes_{ k}\overline{K}$ to  $A_{\overline{K}}$ (cf. \cite[$\mathsection$7]{Lang:AbelianBook} or~\cite[$\mathsection$6]{Conrad}).
If $\Char k=0$, $\tr$ is a closed immersion and $A^{\overline{K}/ k}\otimes_{k}\overline{K}$ can be naturally viewed 
as an abelian subvariety of $A_{\overline{K}}$;
if $\Char k>0$, $\tr$ is a purely inseparable isogeny to its image. 

A {\bf{torsion subvariety}} of $A_{\overline K}$ is a translate $a+C$ of an abelian 
subvariety $C\subset A_{\overline K}$ by a torsion point $a$ of $A_{\overline K}$. A subvariety $X$ of $A_{\overline{K}}$ is said to be {\bf{special}} if  
\begin{equation}
X=\tr(Y{\otimes_{k}\overline{K}})+T
\end{equation}
for some torsion subvariety $T$ of $A_{\overline{K}}$ and  some subvariety $Y$ of $A^{\overline{K}/k}$.
When $X$ is special,  $X(\epsilon)$ is Zariski dense in $X$ for all $\epsilon >0$ (\cite[Theorem 5.4, Chapter 6]{Lang1983}).



%%%%%%%%%%%%%%%%%%%%%%%%%%%%%%%%%%%%%%%%%%%%%%
\subsection{Geometric Bogomolov conjecture}
%%%%%%%%%%%%%%%%%%%%%%%%%%%%%%%%%%%%%%%%%%%%%%

The goal of this paper is to prove the following theorem, which is known as the geometric Bogomolov conjecture. 

\begin{thm}[geometric Bogomolov conjecture]
\source{geometric-bogomolov-conjecture-arbitrary-characteristics}
\notready \label{thmmain} 
\source{geometric-bogomolov-conjecture-arbitrary-characteristics}
Let $k$ be an algebraically closed field.
Let $K/k$ be a finitely generated field extension of transcendence degree at least 1.
Let $A$ be an abelian variety over $K$. 
Let $X$ be a closed subvariety of $A_{\overline{K}}$. 
If $X$ contains a dense set of small points of $A/K/k$, then $X$ is special.
\end{thm}

The above geometric Bogomolov conjecture was proposed by Yamaki \cite[Conjecture 0.3]{Yamaki2013}, but particular
 instances of it were studied earlier by Gubler in~\cite{Gubler2007}. It is an analog over function fields of the original Bogomolov conjecture over number fields which was proved by Ullmo \cite{Ullmo1998} and Zhang \cite{Zhang1998}.

Let us give a quick historical recall on results about the Bogomolov conjecture and its geometric version.
The original Bogomolov conjecture over number fields was proved by Ullmo \cite{Ullmo1998} and Zhang \cite{Zhang1998}, where the major techniques are the equidistribution theorem of Szpiro--Ullmo--Zhang \cite{SUZ}.
The treatments are extended in terms of the Moriwaki height to finitely generated fields over number fields by Moriwaki \cite{Moriwaki2000}.

For the geometric Bogomolov conjecture, it was first proved by Gubler \cite{Gubler2007} when $A$ is totally degenerate at some place of $K/k$.
Yamaki \cite{Yamaki2016c, Yamaki2016a} reduced the geometric Bogomolov conjecture to the case of abelian varieties with good reduction everywhere and with a trivial trace; based on this reduction theorem, Yamaki \cite{Yamaki2017} proved the conjecture for $\dim(X)=1$ or $\mathrm{codim}(X)=1$; he also proved the conjecture for abelian varieties of dimension 5, good reduction everywhere and with trivial trace in \cite{Yamaki2017a}.

The works of Gubler \cite{Gubler2007} and Yamaki \cite{Yamaki2016c} are 
valid in arbitrary characteristics, and extend the strategy of Ullmo and Zhang applying equidistribution theorems. 
In fact, Gubler \cite{Gubler2007} considered tropical varieties of subvarieties of abelian varieties over non-archimedean fields, 
extended the equidistribution theorem of Szpiro--Ullmo--Zhang to the tropical setting, and studied the equilibrium measure in this setting.
Yamaki \cite{Yamaki2016c} did more careful analysis of the situation using equidistribution over Berkovich spaces. The latter equidistribution theorem was proved by Gubler \cite{Guber2008} and Faber \cite{Faber2009b}, which generalized the equidistribution theorems over number fields of Szpiro--Ullmo--Zhang \cite{SUZ},  Chambert-Loir \cite{CL2006}, and Yuan \cite{Yuan2008}.

Before Yamaki \cite{Yamaki2017}, various examples and partial results of the geometric Bogomolov conjecture with $\dim(X)=1$ were previously obtained by \cite{Parshin, Moriwaki1996, Moriwaki1997, Moriwaki1998, Yamaki2002, Yamaki2008, Faber2009, Cinkir2011}.
In particular, Cinkir \cite{Cinkir2011} proved the geometric Bogomolov conjecture for 
$\dim(X)=1$ and $\Char k=0$, based on a height identity of Zhang \cite{Zhang2010}.

In the case $\Char k=0$, Gao--Habegger  \cite{Gao2019} proved the geometric Bogomolov conjecture for $\trdeg(K/k)=1$. 
Recently Cantat--Gao--Habegger--Xie \cite{Cantat2021} proved the full geometric Bogomolov conjecture for $\Char k=0$. 
These proofs are based on the Betti map in the complex-analytic setting.





\subsection{Plan of proof}

Our proof of the geometric Bogomolov conjecture is based on Yamaki's reduction theorem, which reduces the conjecture to the case of good reduction and trivial trace.
We also need the Manin--Mumford conjecture (in the case of trivial trace) proved by Raynaud and Hrushovski.

First, we reduce the conjecture to the case that $K/k$ has transcendence degree 1. This is the main goal of \S\ref{sec trdeg}.
The idea is to take intermediate fields $k'$ of $K/k$, algebraically closed in $K$ and with transcendence degree 1 over $k$, such that the geometric Bogomolov conjecture for $(A/K/k,X)$ follows from those of $(A/K/k',X)$.
For the construction, let $(S,\sM)$ be a polarization of $K/k$.
Use $\sM$ to get a pencil of hyperplane sections of $S$ 
parametrized by $\P_k^1$. Take $k'=k(\P_k^1)$ to be the function field.
There is a generic hyperplane section $H$ over the generic point $\Spec k'$ of $\P_k^1$. 
In particular, $(H,\sM|_H)$ is a polarization of $K/k'$.  
The process from $(S,\sM)$ to $(H,\sM|_H)$ does not increase canonical heights of subvarieties of $A_\OK$. 
Then we carry out careful analysis of the change of special subvarieties of $A$ under this process.

We remark that a well-known procedure to reduce the transcendence degree is to take a closed hyperplane section of 
of $S$ over $k$. This process reduces $K/k$ to $K'/k$ (instead of our $K/k'$), and thus changes $X$ and $A$ by the corresponding reductions. Because of this, it is hard to track density of small points, and it is also hard to track the change of the trace of $A$.
Our method does not change $K$ and thus make everything trackable.

As a consequence, we can assume that $K/k$ has transcendence degree 1. 
Then we apply Yamaki's reduction theorem. 
This reduces the problem to the essential case that $K/k$ has transcendence degree 1, 
and $A$ has good reduction over $S$ and trivial $K/k$-trace. 
Here $S$ is the unique smooth projective curve over $k$ with function field $K$.
Then $A\to \Spec K$ extends to an abelian scheme $\pi:\sA\to S$. 


The second step is to prove the Bogomolov conjecture in the essential case. 
Denote by $\sX$ the Zariski closure of $X$ in $A$.
Let $\sL$ be a symmetric, relatively ample and rigidified line bundle over $\sA$.
Our key property is that any torsion multi-section $\sT$ of $\sA\to S$ is numerically equivalent to a multiple of the self-intersection $\sL^{\dim A}$ in the Chow group of 1-cycles in $\sA$.
This is proved in \S\ref{sec line bundle}.

Now we come to \S\ref{sec final}, which is the core of the proof of the essential case.
To illustrate the idea, assume that there is a positive integer $r$ such that the summation map $f:X^r\to A$ is surjective and generically finite. This happens, for example, if $X$ is a curve and $A$ is the Jacobian variety of $X$.
We have a morphism $f:\sX_{/S}^r\to \sA$ over $S$. 
Note that $X^r\subset A^r$ and $\sX_{/S}^r\subset \sA_{/S}^r$.
Moreover, $X^r$ has canonical height 0 in $A^r$, by Zhang's fundamental inequality and the assumption that $X$ has a dense set of small points. 

Let $\sT$ be a torsion multi-section of $\sA$. 
Let us first assume that $\sT'=f^{-1}(\sT)$ is finite and flat over $S$, and address the technical issue later. 
Denote by $h:\sA_{/S}^r\to \sA$ the summation morphism. 
For any symmetric, relatively ample and rigidified line bundle 
$\sL_r$ over $\sA_{/S}^r$, we have
$$
[\sT']\cdot \sL_r^{\dim \sT'}=[\sX_{/S}^r] \cdot h^*[\sT]\cdot \sL_r^{\dim \sT'}
=a\, [\sX_{/S}^r] \cdot h^*(\sL^{\dim A})\cdot \sL_r^{\dim \sT'}
=0.
$$
Here $a>0$ is a constant coming from the numerical equivalence mentioned above. 
The last equality follows from the fact that $X^r$ has canonical height 0 in $A^r$.
This implies that $T'=\sT'_K$ has canonical height 0, and thus is a finite set of torsion points of $A^r$. 
When varying $\sT$, we obtain a Zariski dense set of torsion points of $X^r$ in $A^r$. 
By the Manin-Mumford conjecture, $X^r$ is torsion, and thus $X$ is torsion. 

In general, we are not able to find $r$ such that $f:X^r\to A$ is surjective and generically finite.
But we can manage to find $r$ such that $f$ is surjective (up to replacing $A$ by an abelian subvariety), and that the relative dimension $e$ of $f$ is strictly smaller than $\dim X$.
If $\sT'=f^{-1}(\sT)$ is flat over $S$ of dimension $e+1$, then the above process still implies that $T'=\sT'_K$ has canonical height 0. 
Then we can conclude that $T'$ is a torsion subvariety in $A^r$ by induction, since $\dim T'<\dim X$.
Torsion points of $T'$ are also torsion points of $X^r$. 
Varying $T'$, we obtain a dense set of torsion points of $X^r$ in $A^r$. 
Then the Manin-Mumford conjecture implies that $X$ is torsion. 


In the above, we have made a technical assumption that $\sT'=f^{-1}(\sT)$ is flat over $S$ of dimension $e+1$.
To remove this caveat, we will have to treat the case that $\sT'=f^{-1}(\sT)$ has irreducible components of dimension bigger than the expected dimension. Note that it is easy to choose $\sT$ so that the Zariski closure $\sT^*$ of $T'=\sT'_K$ in $\sT'$ has the correct dimension. 
Then we prove that, the difference $[\sX_{/S}^r] \cdot h^*[\sT]-[\sT^*]$ is linearly equivalent to an effective Chow cycle of the correct dimension. 
It remedies the above argument. 
This is the content of \S\ref{sec nonproper}, and was inspired by a result of Jia--Shibata--Xie--Zhang.


The idea of converting the Bogomolov conjecture to the Manin--Mumford conjecture was originally used by Zhang \cite{Zhang1992, Zhang1995a} in his proof of the Bogomolov conjecture for powers of $\G_m$. 
In particular, \cite[Lemma 6.6]{Zhang1995a} produced enough torsion hypersurfaces of $X$, and thus enough torsion points of $X$ by induction.  
On the other hand, our situation is more complicated due to lack of torsion hypersufaces in $A$, and our solution is based on the crucial numerical identity between $\sT$ and $\sL^{\dim A}$. As these cycles are not of codimension 1, it also brings issues of non-proper intersection discussed above.



\subsection{Notation and Terminology}
\begin{points}
\item[$\bullet$] For any field $F$, denote by $\overline F$ its algebraic closure.
\item[$\bullet$] For a field extension $K/k$, denote by $\trdeg(K/k)$ the transcendence degree.
\item[$\bullet$] A \emph{variety} is an integral separated scheme of finite type over a field.
\item[$\bullet$] For a Cartier divisor $H$ on a scheme $X$, denote by $|H|$ the linear system associated to $H.$
\item[$\bullet$] For an integral scheme $X$, denote by $\eta_X$ its generic point.
\item[$\bullet$] For a scheme $X$ over a field, denote by $X^{\sm}$ its smooth locus.
\item[$\bullet$] A \emph{subvariety} of a variety is a closed integral subscheme.
%\item[$\bullet$]For a rational map $f: X\dashrightarrow Y$ between varieties. Denote by $I(f)$ the indeterminacy locus of $f$. For a subvariety $V$ of $X$, if $I(f)$ does not contain any irreducible component of $V$, we write $f_{\#}(V)$ for the strict transform $\overline{f(V\setminus I(f))}$ of $V$.
%\item[$\bullet$]For a subscheme $Z$ of a scheme $X$, denote by $I_Z$ the ideal sheaf associated to $Z$.
\item[$\bullet$]
By a \emph{line bundle} over a scheme, we mean an invertible sheaf over the scheme. 
We often write or mention tensor products of line bundles additively, so $aL-bM$ means
$L^{\otimes a}\otimes M^{\otimes (-b)}$
for line bundles $L,M$ and integers $a,b$.
\end{points}


\medskip


\noindent\textbf{Acknowlegement.}
The authors would like to thank the support of the China-Russia Mathematics Center during the preparation of this paper.
The first-named author would like to thank Beijing International Center for Mathematical Research in Peking University 
for the invitation. The second-named author would also like to thank the hospitality of Shandong University, Qingdao for a workshop in July 2021.
